% Source: https://artofproblemsolving.com/community/c3h1428370_number_theory__variables
% Stable identity: aops-topic-1428370-q01
Let $f(x)=x(4-x)$, so $q=f(p)$, $r=f(q)$, $p=f(r)$: the three numbers form a cycle of length $3$ of $f$.

\textbf{All three lie in $[0,4]$.} If $x<0$ then $f(x)=x(4-x)<x$ (as $4-x>1$), and $f(x)<0$ again. If one of the numbers, say $p$, were negative, we would get $p>q>r>p$, impossible. If $p>4$ then $q=f(p)<0$, and the same argument applied to $q$ gives a contradiction. By symmetry all three lie in $[0,4]$.

\textbf{Substitution.} Write $p=2-2\cos\theta$ with $0\le\theta\le\pi$. Since
\[
f(2-2\cos\theta)=(2-2\cos\theta)(2+2\cos\theta)=4\sin^2\theta=2-2\cos2\theta,
\]
we get $q=2-2\cos2\theta$, $r=2-2\cos4\theta$ and $p=2-2\cos8\theta$. So $\cos8\theta=\cos\theta$, that is $8\theta=\pm\theta+2k\pi$, so $\theta$ is a multiple of $\dfrac{2\pi}{7}$ or of $\dfrac{2\pi}{9}$.
\begin{itemize}
\item $\theta=0$ gives $p=q=r=0$, and $\theta=\dfrac{2\pi}{3}=\dfrac{6\pi}{9}$ gives $p=q=r=3$ (fixed points of $f$): excluded, as the numbers must be distinct.
\item $\theta=\dfrac{2\pi}{7}$ (and $\dfrac{4\pi}{7}$, $\dfrac{6\pi}{7}$, which give the same three numbers in another order):
\[
p+q+r=6-2\left(\cos\dfrac{2\pi}{7}+\cos\dfrac{4\pi}{7}+\cos\dfrac{8\pi}{7}\right)=6-2\left(-\dfrac{1}{2}\right)=7 .
\]
\item $\theta=\dfrac{2\pi}{9}$ (and $\dfrac{4\pi}{9}$, $\dfrac{8\pi}{9}$, the same numbers):
\[
p+q+r=6-2\left(\cos40^\circ+\cos80^\circ+\cos160^\circ\right)=6-2\cdot0=6 .
\]
\end{itemize}
Here $\cos\dfrac{2\pi}{7}+\cos\dfrac{4\pi}{7}+\cos\dfrac{6\pi}{7}=-\dfrac{1}{2}$ (the real parts of the nontrivial seventh roots of unity sum to $-1$), with $\cos\dfrac{8\pi}{7}=\cos\dfrac{6\pi}{7}$; and $\cos40^\circ+\cos80^\circ=2\cos60^\circ\cos20^\circ=\cos20^\circ=-\cos160^\circ$.

\begin{asy}
usepackage("amsmath");
size(5.5cm);
defaultpen(fontsize(10pt));
import graph;
real f(real x){return x*(4-x);}
real p=2-2*cos(2pi/7), q=2-2*cos(4pi/7), r=2-2*cos(8pi/7);
draw((0,0)--(4.4,0),Arrow(TeXHead)); draw((0,0)--(0,4.4),Arrow(TeXHead));
draw(graph(f,0,4),blue+linewidth(0.9));
draw((0,0)--(4.2,4.2),gray);
draw((p,p)--(p,q)--(q,q)--(q,r)--(r,r)--(r,p)--(p,p),red+linewidth(0.8));
dot((p,q),red); dot((q,r),red); dot((r,p),red);
draw((p,0)--(p,p),dotted); draw((q,0)--(q,q),dotted); draw((r,0)--(r,p),dotted);
label("$p$",(p,0),S,red); label("$q$",(q,0),S,red); label("$r$",(r,0),S,red);
label("$y=x(4-x)$",(-0.1,3.7),W,blue); label("$y=x$",(4.2,4.2),E,gray);
label("\mbox{}",(min(currentpicture,user=true)+max(currentpicture,user=true))/2,black); // keeps the text after the figure black when LaTeX typesets the labels inline
\end{asy}

(The figure shows the $3$-cycle for $\theta=\dfrac{2\pi}{7}$ on the graph of $f$.) Both cycles consist of three distinct numbers, so the possible values are
\[
p+q+r=\boxed{6\ \text{or}\ 7}.
\]
